\scanpage{161}
\mspage{V.33}
% archive sheet number: 160

6.5. Utilisons maintenant les résultats du § 2, qui \unsure{impliquent|indiquent}, \deleted{\unsure{si} $n=\dim{\unsure{\mathscr{O}|\mathcal{O}}}_{X_0,x}$} pour
\[
\tag*{(6.5.1)}n=\dim\mathscr{O}_{X_0,x}
\]
que l'on a
\[
\tag*{(6.5.2)}H^i(V_{\bar{\unsure{\eta|y}}})\deleted{\unsure{\Lambda|\Lambda_\nu}}=0\quad\text{si }i>n.
\]
\deleted{\unsure{Complétons}}\inserted{Utilisons les} \unsure{deux|des} \deleted{la} suites exactes déduites de \deleted{6.3.4} la suite spectrale de Hochschild-Serre
\[
\deleted{H^i(V)\unsure{-|\longrightarrow}H^i(V_{\bar\eta}}
\]
\[
H^{\unsure{*|t}}(V_\eta)\Longleftarrow E_2^{pq}:H^p(\unsure{\pi|\Pi},H^q(V_{\bar\eta}))
\]
[déduite de $\mathbb{R}\unsure{\Gamma|\mathscr{T}}(V_\eta)\simeq\mathbb{R}\Gamma^\pi\mathbb{R}\Gamma(V_{\bar\eta})$] :
\[
\tag*{(6.5.3)}\unsure{0|:}\unsure{-|\longrightarrow}H^{i\unsure{-|+}1}(V_{\bar\eta})_\pi\longrightarrow H^i(V_\eta)\unsure{-|\longrightarrow}H^i(V_{\bar\eta})^\pi\longrightarrow0,
\]
on trouve donc
\[
\tag*{(6.5.4)}H^i(V_\eta)=0\quad\text{si }i>n+1.
\]
\deleted{Utilisons} Écrivons la suite exacte déduite de 6.3.4
\[
\tag*{(6.5.5)}\longrightarrow H^{i-\unsure{2|1}}(V_s)(-1)\deleted{\unsure{\longrightarrow H^i(V)}}\longrightarrow H^i(V)\longrightarrow H^i(V_\eta)\longrightarrow H^{i-1}(V_s)\longrightarrow H^{i\unsure{+|-}1}(V)\longrightarrow\cdots,
\]
on \deleted{trouve} voit que 6.5.4 s'exprime aussi par \unsure{la relation|les relations} :
\[
\tag*{(6.5.6)}\begin{aligned}
i_*:H^{\deleted{\unsure{i-2}}i}\unsure{(V_s,\Lambda)|(V_s)(-1)}&\longrightarrow H^{i+2}(V_{\deleted{\unsure{s|\eta}}})\quad\text{un isom. si }i\deleted{\unsure{\geq}}>n\deleted{\unsure{-1|+1}}\\
&\qquad\text{un épimorph. si }i=n\deleted{\unsure{-1|+1}}
\end{aligned}
\]

\unsure{6}
